\section{Overview}
\label{sec:overview}

\textsf{Y} is a small imperative language with a static checker, designed
so that a complete implementation of the language fits in one head: the
compiler is written in \textsf{Y} itself (\ic{src/y/}) and compiles a
program to a \emph{static} x86-64 ELF executable using only
its own assembler and ELF writer --- no system assembler, linker, or C
toolchain is involved at run time or at build time of user programs.

Two further targets come from the same pipeline: \ic{--target=win}
writes a PE32+ executable (Section~\ref{sec:toolchain}) and
\ic{--target=kernel} writes a freestanding image that boots on bare metal
(Section~\ref{sec:os}).

\subsection*{Implementations}

The compiler in \ic{src/y/} is self-hosting and byte-reproducible.
A generated C seed (\ic{bootstrap/ycc.c}) lets GCC build it without a Y
executable; \ic{src/c/} contains the seed's boxed value runtime.
Compiler passes are maintained in Y.

This document defines observable behavior, including value sharing,
integer wraparound and trap conditions.

\subsection*{Compiler options}

\ic{ycc [OPTIONS] FILE}, \ic{ycc run [OPTIONS] FILE [-- ARG...]}, and
\ic{ycc build [OPTIONS] [DIR]} share their code-generation options.
\ic{--optimize=speed} (the default) and \ic{--optimize=size} select mutually
exclusive goals; there are no numeric optimization levels. Optimizations that
improve both code size and execution speed run in both modes. These include
removing redundant moves, branches to the following label, and dead instructions,
and using zero idioms where arithmetic flags are dead. Speed mode duplicates
short return sequences to remove taken branches; size mode uses compact
constant loads. These final-link passes also optimize bundled standard-library
units and cached packages. Neither mode changes language semantics.

\ic{--emit=asm|source|wat} selects assembly, expanded source, or a wasm
listing (the executable is the default). \ic{--output=PATH} (or \ic{-o PATH})
chooses an executable or assembly path; assembly and the wasm listing default
to stdout, and expanded source goes to stdout. \ic{--target=linux|win|kernel|web}
names the OS and \ic{--arch=x64|aarch64} the instruction set (Linux only; the
default arch is the host's). \ic{--debug} (or \ic{-g}) includes symbols and
source lines.
\ic{--obfuscate[=MODE]} is an opt-in code-protection mode; it never changes what a
program does, and there is one level: the bare flag and \ic{--obfuscate=full}
select the strongest preset. The program's own
bodies are rewritten at final link: numeric constants are rebuilt, stack frames
are shifted, dead bodies, opaque guards and junk bytes are added, calls may
become indirect, basic blocks are reordered, and the diversity-default static
pass is applied. The image is encrypted and starts through a plaintext page;
code pages stay encrypted until first use (a page fault decodes one page), and
eligible helpers, including float/SSE and division shapes, become generated
bytecode run by a per-build interpreter with a per-function opcode map.
Runtime messages are decoded eagerly because a system call reading a protected
page fails with EFAULT instead of taking the page-fault handler; literals and
the bytecode are decoded page by page, and a literal use touches one byte per
page first. \ic{--obfuscate-seed=N} makes a build reproducible;
\ic{--obfuscate-key=PATH} keeps the stream key in a 16-byte external file;
the sub-flags \ic{--obfuscate-vm=off|auto|all|bytecode} (default \ic{bytecode}),
\ic{--obfuscate-layout=source|shuffle}, \ic{--obfuscate-std=off|auto|all} and
\ic{--obfuscate-diversity=standard|high} tune a layer for tests and bring-up;
\ic{--obfuscate-no-anti-debug} disables the dumpable/tracer checks;
\ic{--obfuscate-report} prints a per-build id and counts;
\ic{--obfuscate-map=PATH} writes a label/address map outside the binary. The
cipher and loader are per-build and page-keyed, with
runtime-dependent opaque branches, jumps over random bytes and a key-derived
startup page schedule; there is no plaintext per-function table. A keyed hash
of the sealed ciphertext refuses a modified file, a keyed per-page MAC refuses
a page whose plaintext is wrong when it first decrypts, and a key-derived
canary replaces the old fixed magic. Full mode reads TracerPid instead of becoming a
tracee, so lazy page faults stay with the loader. Neither makes a standalone
binary cryptographically secure.
Obfuscation covers Linux ELF and Windows PE; a PE has a plaintext entry stub
that decrypts code and literals in place in a writable \ic{.text}, uses the
flatten layer and the embedded key, and has no lazy pages or per-page MACs.
It refuses \ic{-g}, \ic{--pack}, \ic{--emit=asm}, \ic{--emit=source} and \ic{--target=kernel}. The mode raises the cost of static and
dynamic analysis; it is not a cryptographic guarantee.
\ic{--warnings=none|default|extra|all|fix} selects diagnostics by name
(\ic{fix} applies the safe shortenings in place). \ic{--gc=never|auto|always|stress},
\ic{--regions=auto|off}, \ic{--exec=on|off}, and \ic{--std=pack|source} control
the collector, automatic regions, the \ic{exec()} builtin, and standard-unit
sources. \ic{--pack[=METHOD]} writes a self-decompressing executable and
\ic{--time} prints a phase table. Value options accept either
\ic{--name=value} or \ic{--name value}; \ic{-O}, \ic{-o} and \ic{-p} are the
short forms. Build chooses its own output paths and accepts executable output
only.

The default is \ic{--gc=never}: executables contain no collector. The
lifetime solver places temporary allocations in regions and rejects a
program when it cannot prove reclamation for repeated allocations.
\ic{--gc=auto} explicitly permits collection where needed;
\ic{--gc=always} forces it; \ic{--gc=stress} runs a cycle at every allocation.

\subsection*{Notation}

Grammar productions are given in EBNF, with
\texttt{typewriter} for terminals, \ty{italic} for nonterminals,
\texttt{[~]} for option, \texttt{\{~\}} for repetition,
\texttt{|} for alternation, and \texttt{(~)} for grouping.
Semantic rules are written as equational desugars or prose; the desugars
are normative: an implementation may use any mechanism that produces the
same observable behavior.

\subsection*{A taste}

\begin{lstlisting}
// values, references, and the mutating operators
let xs = [1, 2, 3];      // heap array: len 3
xs += 4;                 // in-place push:        [1, 2, 3, 4]
xs *= 0;                 // in-place prepend:     [0, 1, 2, 3, 4]
let ys = xs + [9];       // NEW array, xs untouched

// the target of a compound assignment runs exactly once
let a = [1, 2];
def idx() -> long { print("called"); 0 };
a[idx()] += 10;          // idx() runs ONE time

// records are reference values too
let r = { name: "y", nums: [1, 2] };
r.nums += 3;             // visible through every binding of r
\end{lstlisting}
